%----------------------------------------------------------------------------------------
% TOEPASSINGEN VAN INTEGRALEN. Alle opgaven zijn eigen werk.
%
% Opbouw: Start, leerstof, en op een nieuwe pagina de Opdrachten, Koppel en Test jezelf.
% Labels van dit hoofdstuk beginnen met h4-.
%----------------------------------------------------------------------------------------

\mychapter{4}{Toepassingen van integralen}

\lkr{Leerplan B+S'': de wenken bij LPD 31 en 33. Volume van een omwentelingslichaam en booglengte moet je kunnen \emph{aangeven} (berekeningen eventueel via ICT): hier rekenen we exact waar dat lukt, en verder numeriek met het rekentoestel. De contexten zijn die uit het leerplan: afstand, snelheid, arbeid en energie. Oneigenlijke integralen zijn een mogelijkheid volgens het leerplan (\emph{je kan}): sla ze over als er geen tijd is, want niets anders in de cursus steunt erop, behalve enkele Koppel-oefeningen. Manteloppervlakte staat niet in het leerplan, differentiaalvergelijkingen alleen bij het complementair gedeelte: die ontbreken hier dus.}

Een integraal is meer dan een oppervlakte. Elke grootheid die je kunt opdelen in heel veel kleine stukjes en weer optellen, kun je als integraal schrijven: het volume van een vaas, de lengte van een kabel, de arbeid om een veer uit te rekken, de energie uit zonnepanelen. Dat is de rode draad van dit hoofdstuk. Bij elke toepassing stel je dezelfde vraag: wat is hier één klein stukje, en hoeveel draagt het bij? De som is een Riemannsom, en de limiet is een integraal.

Op het einde gaan we nog een stap verder: wat als het gebied oneindig ver doorloopt? Kan dat een eindige oppervlakte hebben?

\begin{oefeningen}[Start]
\oef{h4-start-1}

\oef{h4-verken}
\end{oefeningen}

\lkr{Oefening \ref{oef:h4-verken} zet de leerlingen zelf op de redenering die in elke toepassing terugkomt: snijd in dunne stukjes, schrijf de bijdrage van een stukje, en tel op. Laat daarna expliciet vragen: wat is hier het stukje?}

%========================================================================================
\section{Volume}

\subsection{Volume uit doorsneden}

Snij een lichaam in dunne plakjes loodrecht op de $x$-as. Een plakje op de plaats $x$ heeft een volume van ongeveer $A(x)\,\Delta x$, met $A(x)$ de oppervlakte van de doorsnede. Optellen en de limiet nemen geeft een integraal.

\begin{theorie}
\begin{stelling}{Volume uit doorsneden}
Is $A(x)$ de oppervlakte van de doorsnede loodrecht op de $x$-as op de plaats $x$, dan is het volume van het lichaam tussen $x = a$ en $x = b$

\vspace{1mm}
\qquad $\formule{V = \int_a^b A(x)\dx}$
\end{stelling}
\end{theorie}

\voorbeeld{}{Toon aan dat een piramide met vierkant grondvlak (zijde $s$) en hoogte $h$ volume $\tfrac13 s^2 h$ heeft.\label{vb:h4-doorsnede}}{%
\uitleg{We leggen de $x$-as door de top, loodrecht op het grondvlak. Dan is elke doorsnede een vierkant, en hangt de zijde af van de plaats $x$.}
\begin{minipage}[t]{.72\linewidth}
\begin{tabular}{@{}>{$\displaystyle}l<{$}@{\hspace{3mm}}>{\RaggedRight\small\itshape\arraybackslash}p{4.6cm}@{}}
\text{zijde op plaats } x = \dfrac{s}{h}\,x & De top ligt in 0. Gelijkvormigheid: bij $x = h$ moet de zijde $s$ zijn, en in het midden de helft. \\[2.5ex]
A(x) = \dfrac{s^2}{h^2}\,x^2 & De oppervlakte van een vierkant is de zijde in het kwadraat. \\[2.5ex]
V = \ig_0^h \dfrac{s^2}{h^2}\,x^2\dx = \dfrac{s^2}{h^2} \cdot \dfrac{h^3}{3} = \dfrac13 s^2 h & $s$ en $h$ zijn constanten: alleen $x$ verandert.
\end{tabular}
\end{minipage}\hfill
\begin{minipage}[t]{.24\linewidth}
\vspace{-2mm}
\begin{tikzpicture}[x=1.3cm,y=1.3cm,font=\scriptsize]
  \draw[lijngrijs,->] (-.1,0) -- (2.2,0) node[right,text=black] {$x$};
  \fill[accent!35] (1.056,.407) -- (1.056,-.165) -- (.814,-.407) -- (.814,.165) -- cycle;
  \draw[accent] (1.056,.407) -- (1.056,-.165) -- (.814,-.407) -- (.814,.165) -- cycle;
  \draw[accent,dashed] (0,0) -- (1.48,.30);
  \draw[accent,thick] (0,0) -- (1.92,.74) -- (1.92,-.30) -- (1.48,-.74) -- cycle;
  \draw[accent,thick] (0,0) -- (1.92,-.30);
  \draw[accent,dashed] (1.48,.30) -- (1.92,.74) (1.48,.30) -- (1.48,-.74);
  \node[above] at (.85,.42) {$A(x)$};
\end{tikzpicture}
\end{minipage}}

\begin{werkwijze}
Kies een $x$-as en de grenzen. Teken een doorsnede op een willekeurige plaats $x$ en druk de oppervlakte uit in $x$, bijvoorbeeld met gelijkvormigheid. Integreer.
\end{werkwijze}

\nujij{Een lichaam ligt tussen $x = 0$ en $x = 2$. De doorsnede op de plaats $x$ is een rechthoek van $x$ bij $2x$. Bereken het volume.}{$A(x) = 2x^2$, dus $V = \tfrac{16}{3}$.}

\begin{oefeningen}
\oef{h4-doorsnede}
\end{oefeningen}

\subsection{Omwentelingslichamen}

Veel voorwerpen zijn draaisymmetrisch: een vaas, een ton, een trechter. Je krijgt ze door een kromme te wentelen om een as. Doorsneden loodrecht op die as zijn dan altijd cirkels, en dat maakt het rekenwerk eenvoudig. Wentel je de grafiek van $f$ om de $x$-as, dan heeft de doorsnede op de plaats $x$ straal $f(x)$, en dus oppervlakte $\pi f(x)^2$. Vul dat in de formule voor doorsneden in.

\begin{theorie}
\begin{stelling}{Omwenteling om de $x$-as}
\begin{minipage}[t]{.60\linewidth}
\vspace{0pt}
Van de grafiek van $f$ op $[a, b]$:

\vspace{1mm}
\qquad $\formule{V = \pi \int_a^b f(x)^2\dx}$
\end{minipage}\hfill
\begin{minipage}[t]{.36\linewidth}
\vspace{0pt}
\begin{tikzpicture}[x=.95cm,y=.72cm,font=\scriptsize]
  \assen{-.1}{4}{-1.9}{2}
  \kromme{.5}{3.3}{1+.4*sin(\x*90)}
  \draw[accent,dashed] plot[domain=.5:3.3,samples=60,smooth] (\x,{-(1+.4*sin(\x*90))});
  \pgfmathsetmacro{\ra}{1+.4*sin(45)}\pgfmathsetmacro{\rb}{1+.4*sin(297)}\pgfmathsetmacro{\rc}{1+.4*sin(180)}
  \draw[accent] (.5,0) ellipse (.13 and \ra);
  \draw[accent] (3.3,0) ellipse (.13 and \rb);
  \fill[accent!30] (2,0) ellipse (.13 and \rc);
  \draw[accent] (2,0) ellipse (.13 and \rc);
  \node[above] at (2,1.05) {$f(x)$};
  \node[below] at (.5,-1.35) {$a$};
  \node[below] at (3.3,-1.0) {$b$};
\end{tikzpicture}
\end{minipage}
\end{stelling}
\end{theorie}

\voorbeeld{}{Toon aan dat een bol met straal $r$ volume $\tfrac43 \pi r^3$ heeft.\label{vb:h4-bol}}{%
\uitleg{Een bol ontstaat door een halve cirkel te wentelen. We schrijven die als functie, en vullen ze in de formule in.}
\naast[]{\begin{stappenvrij}{7.3cm}{3.3cm}
f(x) = \sqrt{r^2 - x^2} \quad \text{op } [-r, r] & De bovenste helft van de cirkel $x^2 + y^2 = r^2$. Wentelen geeft de bol. \\[1ex]
V = \pi \ig_{-r}^{r} (r^2 - x^2)\dx = 2\pi \ig_0^r (r^2 - x^2)\dx & De wortel valt weg door het kwadraat, en de functie is even. \\[2.5ex]
\phantom{V} = 2\pi \Bigl(r^3 - \dfrac{r^3}{3}\Bigr) = \dfrac43 \pi r^3 & Let op: $r$ is hier een constante, en $x$ de veranderlijke.
\end{stappenvrij}}{\begin{tikzpicture}[x=.68cm,y=.68cm,font=\scriptsize]
  \fill[accent!30] (-2,0) arc (180:0:2) -- cycle;
  \assen{-2.4}{2.6}{-2.3}{2.5}
  \draw[accent,thick] (-2,0) arc (180:0:2);
  \draw[accent,dashed] (2,0) arc (0:-180:2);
  \draw[accent,fill=accent!45] (.9,0) ellipse (.22 and 1.786);
  \node[below] at (2,0) {$r$};
\end{tikzpicture}}}

Zo krijgt een formule die je al jaren kent een verklaring. De figuur toont een doorsnede: een cirkelschijf waarvan de straal van plaats tot plaats verandert.

\nujij{Bereken het volume als je $y = x^2$ op $[0, 1]$ wentelt om de $x$-as.}{$\pi \int_0^1 x^4\dx = \tfrac{\pi}{5}$.}

\begin{oefeningen}
\oef{h4-bol}
\end{oefeningen}

Wat als het gebied niet tegen de as ligt? Dan krijgt het lichaam een holte. Kijk eerst naar het eenvoudigste geval, waar je het volume ook zonder integraal vindt.

\voorbeeld{}{Het gebied tussen $y = 2$ en $y = 1$ op $[0, 3]$ wentelt om de $x$-as. Bereken het volume van de buis die zo ontstaat.\label{vb:h4-buis}}{%
\uitleg{Een buis is een grote cilinder waar een kleine cilinder uit gehaald is. We berekenen dus twee volumes, en trekken ze van elkaar af.}
\naast[]{\begin{stappenvrij}{6.1cm}{4.3cm}
V_1 = \pi \ig_0^3 2^2\dx = 12\pi & De grote cilinder: wentel $y = 2$. \\[2.5ex]
V_2 = \pi \ig_0^3 1^2\dx = 3\pi & De kleine cilinder: wentel $y = 1$. \\[2.5ex]
V = V_1 - V_2 = \pi \ig_0^3 (2^2 - 1^2)\dx = 9\pi & Elke doorsnede is een ring: een grote schijf min een kleine.
\end{stappenvrij}}{\begin{tikzpicture}[x=.75cm,y=.75cm,font=\scriptsize]
  \fill[accent!30] (0,1) rectangle (3,2);
  \fill[accent!30] (0,-2) rectangle (3,-1);
  \assen{-.3}{3.7}{-2.3}{2.6}
  \draw[accent,thick] (0,2) -- (3,2) (0,1) -- (3,1);
  \draw[accent,thick,dashed] (0,-2) -- (3,-2) (0,-1) -- (3,-1);
  \draw[accent,even odd rule,fill=accent!45] (3,0) ellipse (.3 and 2) (3,0) ellipse (.15 and 1);
  \node[above left] at (0,2) {2};
  \node[below left] at (0,1) {1};
  \node[below right] at (3,0) {3};
\end{tikzpicture}}}

Dat werkt ook als de randen geen rechten zijn: het lichaam van de bovenste grafiek min het lichaam van de onderste.

\begin{theorie}
\begin{stelling}{Omwenteling tussen twee grafieken}
Het gebied tussen de grafieken van $f$ en $g$ op $[a, b]$, met $f(x) \geqslant g(x) \geqslant 0$: \qquad $\formule{V = \pi \int_a^b \bigl(f(x)^2 - g(x)^2\bigr)\dx}$
\end{stelling}
\end{theorie}

\voorbeeld{}{Het gebied tussen $y = x$ en $y = x^2$ wentelt om de $x$-as. Bereken het volume.\label{vb:h4-ring}}{%
\uitleg{Eerst de snijpunten en de bovenste grafiek. De figuur toont de twee volumes die we van elkaar aftrekken: de kegel van $y = x$, en de holte die $y = x^2$ eruit haalt.}
\figuren{\begin{tikzpicture}[x=2.2cm,y=1cm,font=\scriptsize]
  \lichaam{0}{1}{\x}
  \assen{-.1}{1.3}{-1.25}{1.35}
  \node[above,text=accent] at (.42,.55) {$y = x$};
  \node[below=1pt] at (1.1,0) {1};
\end{tikzpicture}\ \raisebox{12mm}{$-$}\
\begin{tikzpicture}[x=2.2cm,y=1cm,font=\scriptsize]
  \lichaam{0}{1}{(\x)^2}
  \assen{-.1}{1.3}{-1.25}{1.35}
  \node[above,text=accent] at (.42,.25) {$y = x^2$};
  \node[below=1pt] at (1.1,0) {1};
\end{tikzpicture}\ \raisebox{12mm}{$=$}\
\begin{tikzpicture}[x=2.2cm,y=1cm,font=\scriptsize]
  \gebiedtussen{0}{1}{\x}{(\x)^2}
  \gebiedtussen{0}{1}{-(\x)^2}{-\x}
  \assen{-.1}{1.3}{-1.25}{1.35}
  \kromme{0}{1}{\x}
  \kromme{0}{1}{(\x)^2}
  \kromme{0}{1}{-\x}
  \kromme{0}{1}{-(\x)^2}
  \draw[accent] (1,0) ellipse (.14 and 1);
  \draw[accent,even odd rule,fill=accent!45] (.55,0) ellipse (.077 and .55) (.55,0) ellipse (.042 and .3025);
  \node[below=1pt] at (1.1,0) {1};
\end{tikzpicture}}
\begin{stappen}
x = x^2 \ \Rightarrow \ x = 0,\ x = 1 & De snijpunten zijn de grenzen. Op $[0, 1]$ ligt $y = x$ bovenaan, bijvoorbeeld in $x = \tfrac12$. \\[1ex]
V = \pi \ig_0^1 \bigl(x^2 - (x^2)^2\bigr)\dx & Een doorsnede is een ring met buitenstraal $x$ en binnenstraal $x^2$. Beide stralen in het kwadraat. \\[2.5ex]
\phantom{V} = \pi \Bigl(\dfrac13 - \dfrac15\Bigr) = \dfrac{2\pi}{15} & Er blijven enkel machten over: $x^2$ en $x^4$. Dat is de kegel $\tfrac{\pi}{3}$ min de holte $\tfrac{\pi}{5}$.
\end{stappen}}

De klassieke fout: $\pi \int_0^1 (x - x^2)^2\dx$. Die geeft $\tfrac{\pi}{30}$, en dat is iets anders. Trek de kwadraten van elkaar af, en niet de functies.

\nujij{Het gebied tussen $y = 1$ en $y = x^2$ wentelt om de $x$-as. Bereken het volume.}{Op $[-1, 1]$: $\pi \int_{-1}^{1} (1 - x^4)\dx = \tfrac{8\pi}{5}$.}

\begin{oefeningen}
\oef{h4-ring}
\end{oefeningen}

\subsection{Om de verticale as}

Een kom of een trechter staat rechtop: de as is verticaal. De redenering blijft dezelfde, maar nu met horizontale schijven op hoogte $y$, en met straal $x$. Daarom schrijf je $x$ als functie van $y$.


\begin{theorie}
\begin{stelling}{Omwenteling om de $y$-as}
Van de grafiek van $x = h(y)$ op $[c, d]$: \qquad $\formule{V = \pi \int_c^d h(y)^2\,\mathrm{d}y}$
\end{stelling}
\end{theorie}

\voorbeeld{}{De parabool $y = x^2$ met $x \geqslant 0$ en $0 \leqslant y \leqslant 4$ wentelt om de $y$-as. Bereken het volume van de kom.\label{vb:h4-vert}}{%
\uitleg{De as is verticaal, dus we schrijven $x$ als functie van $y$. Een schijf op hoogte $y$ heeft straal $x$.}
\naast[]{\begin{stappenvrij}{6.6cm}{3.7cm}
x = \sqrt{y} & Uit $y = x^2$ en $x \geqslant 0$. Dus $h(y)^2 = y$. \\[1ex]
V = \pi \ig_0^4 y\,\mathrm{d}y = \pi \cdot 8 = 8\pi & De grenzen zijn de hoogtes 0 en 4, niet 0 en 2.
\end{stappenvrij}}{\begin{tikzpicture}[x=.62cm,y=.52cm,font=\scriptsize]
  \draw[->,grijs] (-2.6,0) -- (2.8,0) node[right,text=black] {$x$};
  \draw[->,grijs] (0,-.2) -- (0,4.8) node[above,text=black] {$y$};
  \draw[accent,dashed] plot[domain=-2:0,samples=40,smooth] (\x,{(\x)^2});
  \kromme{0}{2}{(\x)^2}
  \draw[accent] (0,4) ellipse (2 and .28);
  \draw[accent,fill=accent!40] (0,2.2) ellipse (1.483 and .2);
  \node[right] at (1.5,2.2) {$x = \sqrt{y}$};
  \node[left] at (0,4.32) {4};
\end{tikzpicture}}}

\begin{werkwijze}
Schrijf $x$ als een functie van $y$. De grenzen zijn nu waarden van $y$.
\end{werkwijze}

\nujij{De rechte $y = 3x$ met $0 \leqslant y \leqslant 6$ wentelt om de $y$-as. Bereken het volume, en controleer met een kegel.}{$x = \tfrac{y}{3}$, dus $V = \pi \int_0^6 \tfrac{y^2}{9}\,\mathrm{d}y = 8\pi$. Kegel: $r = 2$, $h = 6$, en $\tfrac13 \pi \cdot 4 \cdot 6 = 8\pi$.}

\begin{oefeningen}
\oef{h4-vert}
\end{oefeningen}

%========================================================================================
\section{Booglengte}

Hoe lang is een kabel die in een bocht hangt? Een rechte kun je meten, een kromme niet. Maar een klein stukje van een kromme lijkt op een rechte. Vervang de kromme dus door een gebroken lijn, en maak de stukjes steeds kleiner. De figuur en de tabel tonen dat voor $y = x^2$ op $[0, 1]$.

\figuren{\begin{tikzpicture}[x=2cm,y=2cm,font=\scriptsize]
  \rooster{0}{1}{1}{0}{1}{1}
  \kromme[dashed]{0}{1}{(\x)^2}
  \draw[accent,thick] (0,0) -- (.5,.25) -- (1,1);
  \foreach \p/\q in {0/0,.5/.25,1/1}{\fill[accent] (\p,\q) circle (1.3pt);}
  \node[above,text=accent] at (.3,.7) {$n = 2$};
\end{tikzpicture}\qquad\qquad
\begin{tikzpicture}[x=2cm,y=2cm,font=\scriptsize]
  \rooster{0}{1}{1}{0}{1}{1}
  \kromme[dashed]{0}{1}{(\x)^2}
  \draw[accent,thick] (0,0) -- (.25,.0625) -- (.5,.25) -- (.75,.5625) -- (1,1);
  \foreach \p/\q in {0/0,.25/.0625,.5/.25,.75/.5625,1/1}{\fill[accent] (\p,\q) circle (1.3pt);}
  \node[above,text=accent] at (.3,.7) {$n = 4$};
\end{tikzpicture}\qquad\qquad
\begin{tabular}[b]{@{}cc@{}}
\toprule
$n$ & lengte \\
\midrule
1 & 1,414 \\
2 & 1,460 \\
4 & 1,474 \\
8 & 1,478 \\
\bottomrule
\end{tabular}}

Als $x$ met $\Delta x$ toeneemt, heeft één stukje volgens Pythagoras lengte $\sqrt{\Delta x^2 + \Delta y^2}$. Haal je daar $\Delta x$ uit, dan staat er een Riemannsom.

\begin{theorie}
\begin{stelling}{Booglengte}
\begin{minipage}[t]{.64\linewidth}
\vspace{0pt}
Is $f'$ continu op $[a, b]$, dan is de lengte van de grafiek van $f$ op $[a, b]$

\vspace{1mm}
\qquad $\formule{L = \int_a^b \sqrt{1 + f'(x)^2}\dx}$
\end{minipage}\hfill
\begin{minipage}[t]{.32\linewidth}
\vspace{0pt}
\begin{tikzpicture}[x=1cm,y=1cm,font=\scriptsize]
  \kromme{.2}{2.8}{.15+.28*(\x)^2*.5+.05*\x}
  \coordinate (P) at (1,{.15+.14+.05});
  \coordinate (Q) at (2.2,{.15+.14*4.84+.11});
  \draw[accent,dashed] (P) -- (Q);
  \draw[lijngrijs] (P) -- (Q|-P) -- (Q);
  \fill[accent] (P) circle (1.2pt) (Q) circle (1.2pt);
  \node[below] at ($(P)!.5!(Q|-P)$) {$\Delta x$};
  \node[right] at ($(Q|-P)!.5!(Q)$) {$\Delta y$};
\end{tikzpicture}
\end{minipage}
\end{stelling}

\begin{bewijs}
Verdeel $[a, b]$ in $n$ delen met breedte $\Delta x$, en vervang de kromme door de gebroken lijn door de punten $(x_k, f(x_k))$. Volgens Pythagoras heeft de koorde boven het $k$-de deel lengte
\bewijsregel{\sqrt{\Delta x^2 + \Delta y^2} = \sqrt{1 + \Bigl(\dfrac{\Delta y}{\Delta x}\Bigr)^{\!2}}\;\Delta x = \sqrt{1 + f'(c_k)^2}\;\Delta x ,}
met $c_k$ het punt uit de stelling van Lagrange, waarvoor $\tfrac{\Delta y}{\Delta x} = f'(c_k)$. De figuur toont dat: ergens op dat deel loopt de raaklijn evenwijdig met de koorde.
\figuren{\begin{tikzpicture}[x=2cm,y=1.35cm,font=\scriptsize]
  \draw[lijngrijs] (.6,.294) -- (2.2,.294) -- (2.2,2.086);
  \draw[accent,dashed] (.6,.294) -- (2.2,2.086);
  \draw[grijs,thick] (.8,.262) -- (2,1.606);
  \kromme{.4}{2.4}{.15+.4*(\x)^2}
  \fill[accent] (.6,.294) circle (1.2pt) (2.2,2.086) circle (1.2pt);
  \fill[black] (1.4,.934) circle (1.2pt);
  \draw[lijngrijs,dashed] (1.4,.934) -- (1.4,.294);
  \node[below] at (.6,.294) {$x_{k-1}$};
  \node[below] at (2.2,.294) {$x_k$};
  \node[below] at (1.4,.294) {$c_k$};
  \node[right] at (2.2,1.2) {$\Delta y$};
  \node[above left,text=accent] at (1.25,1.05) {koorde};
  \node[text=grijs] at (1.83,.95) {raaklijn};
\end{tikzpicture}}
De lengte van de gebroken lijn is dus een Riemannsom, en voor $n \to +\infty$ volgt
\bewijsregel{L = \lim_{n \to +\infty} \sum_{k=1}^{n} \sqrt{1 + f'(c_k)^2}\;\Delta x = \ig_a^b \sqrt{1 + f'(x)^2}\dx .}
\end{bewijs}
\end{theorie}

\voorbeeld{}{Bereken de lengte van de grafiek van $f(x) = \tfrac23 x^{3/2}$ op $[0, 3]$.\label{vb:h4-boog}}{%
\uitleg{Eerst de afgeleide en het kwadraat ervan. Of je exact kunt verder rekenen, zie je pas daarna: hier valt de wortel uit de integrand weg.}
\naast[]{\begin{stappenvrij}{7cm}{4cm}
f'(x) = \sqrt{x}, \quad 1 + f'(x)^2 = 1 + x & Het kwadraat haalt de wortel weg. \\[1ex]
L = \ig_0^3 \sqrt{1 + x}\dx = \grens{\tfrac23 (1 + x)^{3/2}}{0}{3} = \dfrac{14}{3} & Controle: de koorde is $\sqrt{21} \approx 4{,}58$, en de boog moet langer zijn.
\end{stappenvrij}}{\begin{tikzpicture}[x=.7cm,y=.6cm,font=\scriptsize]
  \rooster{0}{1}{3}{0}{1}{4}
  \kromme{0.01}{3}{(2/3)*(\x)^(1.5)}
  \draw[accent,dashed] (0,0) -- (3,3.464);
  \node[above left,text=accent] at (1.6,1.9) {koorde};
\end{tikzpicture}}}

Die functie is met opzet zo gekozen dat de wortel wegvalt. Voor bijna elke andere functie lukt dat niet.

\voorbeeld{}{Bereken de lengte van de grafiek van $g(x) = x^2$ op $[0, 1]$.\label{vb:h4-boog2}}{%
\uitleg{De functie lijkt eenvoudiger dan de vorige, maar hier blijft de wortel staan. Dan geeft het rekentoestel een antwoord.}
\begin{stappen}
g'(x) = 2x, \quad L = \ig_0^1 \sqrt{1 + 4x^2}\dx \approx 1{,}479 & Geen van onze methodes lukt: numeriek. De koorde is $\sqrt{2} \approx 1{,}414$, en dat past bij de tabel hierboven.
\end{stappen}}

\begin{werkwijze}
Bereken $f'(x)$ en werk $1 + f'(x)^2$ uit. Lukt het om er een volkomen kwadraat in te zien, of een eenvoudige vorm, dan kun je exact rekenen. Lukt dat niet, gebruik dan de numerieke integraal van je rekentoestel. Dat is bij booglengte eerder de regel dan de uitzondering.
\end{werkwijze}

\nujij{Bereken de lengte van de rechte $y = 2x$ op $[0, 3]$ met de formule, en controleer met Pythagoras.}{$\int_0^3 \sqrt{5}\dx = 3\sqrt{5}$. Pythagoras: $\sqrt{3^2 + 6^2} = 3\sqrt{5}$.}

\begin{oefeningen}
\oef{h4-boog}
\end{oefeningen}

%========================================================================================
\section{Contexten}

Het is altijd dezelfde redenering: als de grootheid die je zoekt een product is (snelheid maal tijd, kracht maal afstand, vermogen maal tijd), maar een van de factoren verandert onderweg, dan is ze een integraal. In een klein stukje is die factor bijna constant, en daar geldt het product wel. We bekijken de contexten een voor een, telkens vanuit een voorbeeld.

\subsection{Verplaatsing en afgelegde weg}

In hoofdstuk 1 las je al uit een grafiek hoe ver een karretje reed. Nu kun je ook rekenen als de snelheid geen rechte lijn is.

\voorbeeld{}{Een voorwerp heeft snelheid $v(t) = 4 - t^2$ (in m/s) op $[0, 3]$. Bereken de verplaatsing en de afgelegde weg.\label{vb:h4-weg}}{%
\uitleg{De snelheid wisselt van teken in $t = 2$: tot dan beweegt het voorwerp vooruit, daarna achteruit. Voor de afgelegde weg moeten we daar splitsen.}
\naast[]{\begin{stappenvrij}{6.7cm}{4cm}
\ig_0^3 (4 - t^2)\dt = \grens{4t - \dfrac{t^3}{3}}{0}{3} = 3 \text{ m} & De verplaatsing: de meters achteruit tellen negatief. \\[2.5ex]
\ig_0^2 (4 - t^2)\dt = \dfrac{16}{3} & Het stuk vooruit. \\[2.5ex]
\ig_2^3 (4 - t^2)\dt = -\dfrac73 & Het stuk achteruit. \\[2.5ex]
\dfrac{16}{3} + \dfrac73 = \dfrac{23}{3} \text{ m} & De afgelegde weg: elke meter telt positief.
\end{stappenvrij}}{\begin{tikzpicture}[x=.75cm,y=.27cm,font=\scriptsize]
  \gebied{0}{2}{4-(\x)^2}
  \gebied[accent!15]{2}{3}{4-(\x)^2}
  \draw[black!14,very thin,xstep=1,ystep=2] (0,-6) grid (3,4);
  \draw[->,grijs] (0,0) -- (3.5,0) node[right,text=black] {$t$};
  \draw[->,grijs] (0,-6) -- (0,5) node[above,text=black] {$v$};
  \kromme{0}{3}{4-(\x)^2}
  \foreach \t in {1,2,3} \node[above] at (\t,0) {\t};
  \foreach \t in {-4,4} \node[left] at (0,\t) {\t};
\end{tikzpicture}}}

\begin{theorie}
Bij een snelheid $v$ op $[t_1, t_2]$ is de

\vspace{1mm}
\qquad \textbf{\textcolor{accent}{verplaatsing}} \quad $\formule{\int_{t_1}^{t_2} v(t)\dt}$ \qquad en de \textbf{\textcolor{accent}{afgelegde weg}} \quad $\formule{\int_{t_1}^{t_2} |v(t)|\dt}$
\end{theorie}

\subsection{Snelheid uit de versnelling}

Bij een beweging werk je stap voor stap terug. Uit de versnelling haal je de snelheid, en uit de snelheid de plaats: telkens een primitieve, en telkens is er een beginwaarde nodig.

\voorbeeld{}{Een voorwerp start op $t = 0$ met snelheid 3 m/s en versnelling $a(t) = 2t + 1$ (in m/s²). Hoe ver is het na 3 s?\label{vb:h4-versnelling}}{%
\uitleg{Uit de versnelling halen we de snelheid, en uit de snelheid de afstand. De beginsnelheid bepaalt de constante.}
\naast[]{\begin{stappenvrij}{6.8cm}{4cm}
v(t) = t^2 + t + 3 & Een primitieve van $a$, met $v(0) = 3$. \\[1ex]
\ig_0^3 (t^2 + t + 3)\dt = 9 + \dfrac92 + 9 = \dfrac{45}{2} \text{ m} & De snelheid is positief, dus verplaatsing = afgelegde weg.
\end{stappenvrij}}{\begin{tikzpicture}[x=.85cm,y=.13cm,font=\scriptsize]
  \gebied{0}{3}{(\x)^2+\x+3}
  \draw[black!14,very thin,xstep=1,ystep=5] (0,0) grid (3,15);
  \draw[->,grijs] (0,0) -- (3.4,0) node[right,text=black] {$t$};
  \draw[->,grijs] (0,0) -- (0,17) node[above,text=black] {$v$};
  \kromme{0}{3}{(\x)^2+\x+3}
  \foreach \t in {1,2,3} \node[below] at (\t,0) {\t};
  \foreach \t in {5,10,15} \node[left] at (0,\t) {\t};
\end{tikzpicture}}}

\begin{theorie}
De \textbf{\textcolor{accent}{verandering van de snelheid}} volgt uit de versnelling $a$: \quad $\formule{v(t_2) - v(t_1) = \int_{t_1}^{t_2} a(t)\dt}$
\end{theorie}

\nujij{Een rem geeft een versnelling $a(t) = -6t$ (in m/s²). De startsnelheid is 12 m/s. Wanneer staat het voorwerp stil, en hoe ver is het dan?}{$v(t) = 12 - 3t^2$, dus stilstand voor $t = 2$. Afstand: $\int_0^2 (12 - 3t^2)\dt = 24 - 8 = 16$ m.}

\begin{oefeningen}
\oef{h4-context}
\end{oefeningen}

\subsection{Arbeid}

Bij een constante kracht is de arbeid kracht maal afstand. Verandert de kracht onderweg, dan geldt dat alleen nog voor een klein stukje afstand.

\voorbeeld{}{Een ketting van 10 m en 2 kg per meter hangt van een dak. Je trekt ze helemaal omhoog. Hoeveel arbeid kost dat? Neem $g = 9{,}81$ m/s².\label{vb:h4-arbeid}}{%
\uitleg{De kracht is niet constant, want de ketting die je al omhoog trok, hangt niet meer aan je handen. Daarom schrijven we de kracht als functie van de afstand $x$.}
\naast[]{\begin{stappenvrij}{5.6cm}{4cm}
F(x) = 2 \cdot (10 - x) \cdot 9{,}81 & Na $x$ meter hangt er nog $10 - x$ m, en elke meter weegt 2 kg. De kracht neemt dus af. \\[1ex]
W = \ig_0^{10} 19{,}62\,(10 - x)\dx & Kracht maal afstand, met een veranderlijke kracht. \\[2.5ex]
\phantom{W} = 19{,}62 \cdot 50 = 981 \text{ J} & De integraal van $10 - x$ is een driehoek: 50.
\end{stappenvrij}}{\begin{tikzpicture}[x=.33cm,y=.0165cm,font=\scriptsize]
  \fill[accent!30] (0,0) -- (0,196.2) -- (10,0) -- cycle;
  \draw[black!14,very thin,xstep=2,ystep=50] (0,0) grid (10,200);
  \draw[->,grijs] (0,0) -- (11.2,0) node[right,text=black] {$x$};
  \draw[->,grijs] (0,0) -- (0,225) node[above,text=black] {$F$};
  \draw[accent,thick] (0,196.2) -- (10,0);
  \foreach \t in {2,4,...,10} \node[below] at (\t,0) {\t};
  \foreach \t in {100,200} \node[left] at (0,\t) {\t};
\end{tikzpicture}}}

\begin{theorie}
De \textbf{\textcolor{accent}{arbeid}} door een kracht $F$ op een rechte, van $x_1$ naar $x_2$: \quad $\formule{W = \int_{x_1}^{x_2} F(x)\dx}$
\end{theorie}

\begin{oefeningen}
\oef{h4-arbeid}
\end{oefeningen}

\subsection{Energie}

Een toestel met een constant vermogen verbruikt vermogen maal tijd aan energie. Verandert het vermogen, dan tel je weer kleine stukjes op.

\voorbeeld{}{Een toestel heeft een vermogen $P(t) = 2t$ (in kW) op $[0, 3]$ (met $t$ in uur). Hoeveel energie gebruikt het?\label{vb:h4-energie}}{%
\uitleg{Vermogen maal tijd geeft energie, en het vermogen verandert: dus een integraal. Kilowatt maal uur geeft kilowattuur.}
\naast[]{\begin{stappenvrij}{6.8cm}{4cm}
E = \ig_0^3 2t\dt = \grens{t^2}{0}{3} = 9 \text{ kWh} & De oppervlakte onder de grafiek is ook de driehoek $\tfrac12 \cdot 3 \cdot 6$.
\end{stappenvrij}}{\begin{tikzpicture}[x=.8cm,y=.3cm,font=\scriptsize]
  \fill[accent!30] (0,0) -- (3,6) -- (3,0) -- cycle;
  \draw[black!14,very thin,xstep=1,ystep=2] (0,0) grid (3,6);
  \draw[->,grijs] (0,0) -- (3.4,0) node[right,text=black] {$t$};
  \draw[->,grijs] (0,0) -- (0,7) node[above,text=black] {$P$};
  \draw[accent,thick] (0,0) -- (3,6);
  \foreach \t in {1,2,3} \node[below] at (\t,0) {\t};
  \foreach \t in {2,4,6} \node[left] at (0,\t) {\t};
\end{tikzpicture}}}

\begin{theorie}
De \textbf{\textcolor{accent}{energie}} uit het vermogen $P$: \quad $\formule{E = \int_{t_1}^{t_2} P(t)\dt}$
\end{theorie}

\begin{werkwijze}
Controleer de eenheden: de eenheid van de integraal is de eenheid van de functie maal de eenheid op de as. Neem de beginwaarde mee als er een gegeven is. Maak bij een lichaam dat niet overal hetzelfde beweegt plakjes, en zoek de bijdrage van één plakje.
\end{werkwijze}

%========================================================================================
\section{Oneigenlijke integralen}

Tot hier was de functie continu op een begrensd interval. Maar wat als het gebied oneindig ver doorloopt, of oneindig hoog? In de figuur bij het kader hieronder loopt het gebied onder de grafiek naar rechts en stopt nooit. De hoofdstelling geldt dan niet rechtstreeks: $+\infty$ is geen getal om in te vullen. Toch is er een uitweg, en het idee kennen we al van limieten. Stop het gebied bij een grens $b$, bereken die integraal, en laat $b$ groter worden.

\begin{theorie}
\begin{minipage}[t]{.65\linewidth}
\textbf{\textcolor{accent}{Onbegrensd interval.}} Is $f$ continu op $[a, +\infty[$, dan

\vspace{1mm}
\qquad $\formule{\int_a^{+\infty} f(x)\dx = \lim_{b \to +\infty} \int_a^b f(x)\dx}$

\vspace{1.5mm}
\textbf{\textcolor{accent}{Onbegrensde functie.}} Is $f$ continu op $]a, b]$ met een verticale asymptoot in $x = a$, dan

\vspace{1mm}
\qquad $\formule{\int_a^b f(x)\dx = \limr{c}{a} \int_c^b f(x)\dx}$
\end{minipage}\hfill
\begin{minipage}[t]{.31\linewidth}
\vspace{-1mm}
\begin{tikzpicture}[x=.9cm,y=.85cm,font=\scriptsize]
  \gebied{1}{3.6}{1.6/(\x)^2}
  \assen{-.1}{4}{-.1}{2.2}
  \kromme{.75}{3.7}{1.6/(\x)^2}
  \node[below] at (1,0) {$a$};
  \node at (3.45,.4) {$\to$};
\end{tikzpicture}
\end{minipage}

\vspace{1.5mm}
Is die limiet een getal, dan \textbf{\textcolor{accent}{convergeert}} de integraal. Anders \textbf{\textcolor{accent}{divergeert}} ze.

\vspace{1mm}
Hetzelfde geldt bij een grens $-\infty$ en bij een asymptoot in $x = b$. Zit het probleem in de twee grenzen of ergens tussen de grenzen, splits dan in stukken met elk één probleem. De integraal convergeert alleen als elk stuk convergeert.
\end{theorie}

\voorbeeld{}{Bereken, als dat kan, $\ig_1^{+\infty} \dfrac{1}{x^2}\dx$.\label{vb:h4-oneindig}}{%
\uitleg{We vervangen $+\infty$ door een letter $b$, rekenen gewoon, en nemen pas op het einde de limiet.}
\naast[]{\begin{stappenvrij}{6.6cm}{3.2cm}
\lim_{b \to +\infty} \grens{-\dfrac1x}{1}{b} = \lim_{b \to +\infty} \Bigl(1 - \dfrac1b\Bigr) = 1 & Convergent. Een oneindig gebied kan een eindige oppervlakte hebben.
\end{stappenvrij}}{\begin{tikzpicture}[x=.68cm,y=.9cm,font=\scriptsize]
  \gebied{1}{6}{1/(\x)^2}
  \rooster{0}{1}{6}{0}{1}{2}
  \kromme{.71}{6}{1/(\x)^2}
\end{tikzpicture}}}

\voorbeeld{}{Bereken, als dat kan, $\ig_1^{+\infty} \dfrac{1}{x}\dx$.\label{vb:h4-oneindig2}}{%
\uitleg{Dezelfde aanpak, bij een functie die er sterk op lijkt. Toch is de uitkomst heel anders: deze functie daalt te traag.}
\naast[]{\begin{stappenvrij}{6.6cm}{3.2cm}
\lim_{b \to +\infty} \grens{\ln x}{1}{b} = \lim_{b \to +\infty} \ln b = +\infty & Divergent, al lijken de twee functies sterk op elkaar.
\end{stappenvrij}}{\begin{tikzpicture}[x=.68cm,y=.9cm,font=\scriptsize]
  \gebied{1}{6}{1/\x}
  \rooster{0}{1}{6}{0}{1}{2}
  \kromme{.5}{6}{1/\x}
\end{tikzpicture}}}

Hetzelfde idee werkt als de functie zelf naar oneindig gaat. Dan stop je net voor het probleem, en laat je die grens naderen. Let dus niet alleen op de grenzen: de asymptoot kan ook midden in het interval liggen.

\voorbeeld{}{Bereken, als dat kan, $\ig_0^4 \dfrac{1}{\sqrt{x}}\dx$.\label{vb:h4-asymptoot}}{%
\uitleg{De functie gaat naar oneindig in $x = 0$, en 0 is een grens. We stoppen dus net rechts van 0, in $c$, en laten $c$ naderen.}
\naast[]{\begin{stappenvrij}{7cm}{3.8cm}
\limr{c}{0} \grens{2\sqrt{x}}{c}{4} = \limr{c}{0} (4 - 2\sqrt{c}) = 4 & Convergent. De $>$ onder de limiet zegt dat $c$ van rechts nadert: $c > 0$.
\end{stappenvrij}}{\begin{tikzpicture}[x=.5cm,y=.42cm,font=\scriptsize]
  \gebied{.0625}{4}{1/sqrt(\x)}
  \rooster{0}{1}{4}{0}{1}{4}
  \kromme{.0625}{4}{1/sqrt(\x)}
\end{tikzpicture}}}

\voorbeeld{}{Bereken, als dat kan, $\ig_{-1}^{1} \dfrac{1}{x^2}\dx$.\label{vb:h4-midden}}{%
\uitleg{Ook hier gaat de functie naar oneindig in $x = 0$, maar nu ligt 0 binnen het interval. Dat zie je niet aan de grenzen: kijk dus altijd naar het domein.}
\naast[]{\begin{stappenvrij}{7cm}{3.4cm}
\ig_{-1}^{1} \dfrac{1}{x^2}\dx = \ig_{-1}^{0} \dfrac{1}{x^2}\dx + \ig_{0}^{1} \dfrac{1}{x^2}\dx & De functie bestaat niet in 0: splitsen. Dit is de integraal uit hoofdstuk 2. \\[3ex]
\limr{c}{0} \grens{-\dfrac1x}{c}{1} = \limr{c}{0} \Bigl(-1 + \dfrac1c\Bigr) = +\infty & Eén stuk divergeert, dus de hele integraal divergeert.
\end{stappenvrij}}{\begin{tikzpicture}[x=.6cm,y=.42cm,font=\scriptsize]
  \gebied{-1}{-.5}{1/(\x)^2}
  \gebied{.5}{1}{1/(\x)^2}
  \rooster{-2}{1}{2}{0}{1}{4}
  \kromme{-2}{-.5}{1/(\x)^2}
  \kromme{.5}{2}{1/(\x)^2}
\end{tikzpicture}}}

\begin{werkwijze}
Kijk eerst of de integraal oneigenlijk is: een oneindige grens, of een punt van $[a, b]$ dat niet in het domein zit. Schrijf ze als limiet. Bereken de integraal met de letter als grens en neem dan de limiet.
\end{werkwijze}

\nujij{Bereken, als dat kan, $\int_0^{+\infty} \e^{-x}\dx$ en $\int_0^1 \dfrac{1}{x}\dx$.}{$\lim_{b \to +\infty}(1 - \e^{-b}) = 1$; en $\limr{c}{0}(-\ln c) = +\infty$, dus divergent.}

\voorbeeld{}{De andere gevallen uit het kader. (Niet uitgewerkt.)\label{vb:h4-andere}}{%
\uitleg{Schrijf telkens eerst de limiet op: dan is de rest gewoon rekenwerk.}
\begin{tabular}{@{}>{$\displaystyle}l<{$}@{\hspace{6mm}}>{\RaggedRight\small\itshape\arraybackslash}p{8.2cm}@{}}
\ig_{-\infty}^{0} \e^x\dx = 1 & Grens $-\infty$: $\lim_{a \to -\infty}(1 - \e^a) = 1$. \\[2.5ex]
\ig_0^1 \dfrac{1}{\sqrt{1 - x}}\dx = 2 & Asymptoot in de bovengrens: $\liml{c}{1}\bigl(2 - 2\sqrt{1 - c}\bigr) = 2$. \\[3.5ex]
\ig_{-\infty}^{+\infty} \e^{-|x|}\dx = 2 & Twee oneindige grenzen: splits in 0, en elk stuk geeft 1.
\end{tabular}}

\begin{oefeningen}
\needspace{9\baselineskip}
\oef{h4-oneindig}
\end{oefeningen}

%========================================================================================
\opdrachtenkop

\begin{oefeningen}
\oef{h4-o1}

\oef{h4-o3}

\oef{h4-o4}

\oef{h4-o5}

\oef{h4-o6}

\oef{h4-o7}

\oef{h4-o8}

\oef{h4-o9}

\oef{h4-o10}

\oef{h4-o11}

\oef{h4-o12}

\oef{h4-okies}

\oef{h4-o13}

\oef{h4-o14}

\oef{h4-o15}

\oef{h4-o16}

\oef{h4-o17}

\oef{h4-o18}

\oef{h4-o19}

\oef{h4-o21}

\oef{h4-o22}

\oef{h4-o23}

\oef{h4-o24}

\oef{h4-machten}

\oef{h4-o25}

\oef{h4-o26}
\end{oefeningen}

%----------------------------------------------------------------------------------------
% KOPPEL: verbanden met andere leerstof
%----------------------------------------------------------------------------------------
\needspace{20\baselineskip}
\opdrachtsectie{Koppel}

\lkr{Volgens de README van het pakket komen telproblemen (Kansrekenen) vóór Analyse 3, en Rijen, Statistiek en Ruimtemeetkunde erna: kijk na wat er al gezien is. De oefening over Ruimtemeetkunde vraagt enkel de volumeformules van de 2de graad. De dichtheid in de oefening over Statistiek is een voorbeeld; de normale verdeling zelf komt in de cursus Statistiek. De oefeningen over Telproblemen, Statistiek en Rijen vragen de oneigenlijke integralen uit dit hoofdstuk.}

\begin{oefeningen}
\oef{h4-koppel-1}

\oef{h4-koppel-2}

\oef{h4-koppel-3}

\oef{h4-koppel-4}
\end{oefeningen}

%----------------------------------------------------------------------------------------
\opdrachtsectie{Test jezelf}

\begin{oefeningen}
\oef{h4-t1}

\oef{h4-t2}

\oef{h4-t3}

\oef{h4-t4}

\oef{h4-t5}

\oef{h4-t6}

\oef{h4-t7}

\oef{h4-t8}

\end{oefeningen}

\herhaalschema{%
  \fv{oef:h4-t1}{vb:h4-doorsnede}{oef:h4-doorsnede}%
  \fv{oef:h4-t2}{vb:h4-bol}{oef:h4-bol}%
  \fv{oef:h4-t3}{vb:h4-ring}{oef:h4-ring}%
  \fv{oef:h4-t4}{vb:h4-vert}{oef:h4-vert}%
  \fv{oef:h4-t5}{vb:h4-boog}{oef:h4-boog}%
  \fv{oef:h4-t6}{vb:h4-arbeid}{oef:h4-arbeid}%
  \fv{oef:h4-t7}{vb:h4-versnelling}{oef:h4-context}%
  \fv{oef:h4-t8}{vb:h4-asymptoot}{oef:h4-oneindig}}

\ifzelfstudie
\section*{Maak je eigen samenvatting}
Sluit dit boek en schrijf in je schrift: de redenering met plakjes, de formules voor volume met doorsneden en voor omwenteling om elke as, de formule voor booglengte, de integralen uit de contexten, en wanneer en hoe je een oneigenlijke integraal uitrekent. Vergelijk daarna met de kaders.
\fi
